

\section{Partitionnement par intervalle}

\subsection{Indexation distribuée}


\begin{frame}
\frametitle{Appliquons ces principes à la distribution des données}

A quelques variantes près, on procède comme pour un index non-dense.

\begin{itemize}
  \item La collection est triée sur la clé.
   
        \smallskip
   
       Au moins partiellement (cf. plus loin)
       
  \item On la découpe en fragments d'une taille maximale pré-déterminée (de 64MO à qq GO)
   
   \smallskip
   
        Chaque frament couvre donc un intervalle $[min_f, max_f[$
        
   \item Les fragments sont distribués équitablement sur les serveurs

    \smallskip
    
         Un ou plusieurs fragments par serveur.

   \item Un index sur les fragments est maintenu par le routeur.
 
    \smallskip
   
      \alert{Souvenez-vous\,: l'index est tout petit par rapport à la collection}
      
\end{itemize} 

\alert{Le routeur reçoit les requêtes et détermine le(s) serveur(s) à solliciter.}

\end{frame}

\begin{frame}
\frametitle{L'architecture, avec partionnement par intervalle}

Le routeur connaît les intervalles, et les serveurs associés.

\figSlide{partition-range}

\begin{remark}
\begin{itemize}
   \item Chaque serveur maintient un index local sur ses fragments.
   \item Les fragments sont autant que possible conservés en mémoire.
  \item Parmi les variantes\,: les fragments sur les serveurs peuvent être triés ou pas. 
\end{itemize}
\end{remark}
\end{frame}


\begin{frame}
\frametitle{Dynamicité: comment évolue le système\,?}

L'opération de base est le \alert{split}, division d'un gros fragment en deux petits.

\figSlide{partition-split}

\alert{L'équilibrage} consiste à répartir équitablement les fragments.

\end{frame}

\begin{frame}
\frametitle{Les opérations du routeur}

% A two-columns presentation
\begin{tabular}{l|p{6cm}}
   \parbox{5cm}{
$get(k)$: chercher l'intervalle qui contient $k$, transmettre
au serveur correspondant.

\smallskip

$put(k, d)$: identifier le serveur avec $k$, transmettre l'insertion
 
 \smallskip
 
 $delete(k)$: idem
 
 \smallskip
 
 $range(k_1, k_2)$: transmettre à tous les serveurs gérant un intervalle chevauchant $[k_1, k_2]$
 
 \smallskip
 
 Toute autre opération\,: transmettre à tous les serveurs.
   }  % Inclusion of an XML file
  & 
  \parbox{6cm}{
 
\begin{tabular}{r|l}
Intervalle  & Serveur\\
\hline
$]-\infty, a]$ & A  \\
$[a+1, b]$ & B  \\
$[b+1, c]$ & C  \\
$[c+1, d]$ & B  \\
$[d+1, \infty[$ & A  \\
\hline
\end{tabular}

}

\end{tabular}

\end{frame}


\section{Etude de cas 2\,: \mongodb/}

 
\begin{frame}
\frametitle{Le \textit{sharding} de \mongodb/}

On crée des \textit{replica set} qui vont chacun stocker un sous-ensemble des fragments.

\vfill

Un serveur spécial, le \textit{config server} (répliqué), se charge de stocker la table de routage.

\vfill

Chaque application discute avec un processus routeur, \textit{mongos} 

\vfill

\figSlideWithSize{mongosharding}{8cm}


\end{frame}



\begin{frame}[fragile]
\frametitle{Allons-y pour un petit test}

On lance un \textit{config server}

\begin{smallverbatim}
mkdir /data/configdb
mongod --configsvr --dbpath /data/configdb --port 27019
\end{smallverbatim}

On lance un routeur en lui indiquant où se trouve(nt) le(s) \textit{config server}(s)

\begin{smallverbatim}
mongos --configdb localhost:27019
\end{smallverbatim}

Et deux serveurs (il faudrait des \textit{replica sets} complets en production)
\begin{smallverbatim}
mongod --dbpath /data/node1 --port=30001
mongod --dbpath /data/node3 --port=30002
\end{smallverbatim}

A partir de \textit{mongos}, déclarons les deux serveurs comme \textit{shards}.

\begin{smallverbatim}
mongo --port 27017
sh.addShard("localhost:30000")
sh.addShard("localhost:30001")
sh.enableSharding("nfe204")
\end{smallverbatim}

\end{frame}

\begin{frame}[fragile]
\frametitle{Choisir la clé de partitionnement}

On indique la clé de partitionnement avec la commande suivante\,:

\begin{smallverbatim}
sh.shardCollection("nfe204.videos", { "title": 1 } )
\end{smallverbatim}

Attention au choix de la clé

\begin{itemize}
  \item Les clés-séquences\,:  on envoie toujours les insertions au même serveur\,; bon ou pas?
  \item Les clés à faible cardinalité (ex: pays)\,: pas très bon non plus, car 
               il pourra être difficile de faire un \textit{split}
\end{itemize}

Bonne option: appliquer un hachage pour bien distribuer les insertions\,:

\begin{smallverbatim}
sh.shardCollection("nfe204.videos", { "title": "hashed" } )
\end{smallverbatim}

\end{frame}


\begin{frame}[fragile]
\frametitle{Inspecter la configuration}

Liste des \textit{shards}
\begin{smallverbatim}
db.runCommand({listshards: 1})
\end{smallverbatim}

Statut général du cluster
\begin{smallverbatim}
sh.status()
\end{smallverbatim}

Les fragments sont dans une collection \textit{chuncks}.

\begin{smallverbatim}
use config
db.chuncks.count()
db.chuncks.findOne()
\end{smallverbatim}
\end{frame}


\section{Etude de cas 1\,: BigTable (HBase, Cassandra, ...)}

\begin{frame}
\frametitle{BigTable\,: la structure} 

C'est une collection \alert{triée}, dont chaque ligne est indexée par une clé.


 \figSlideWithSize{bigtable}{8cm}
 
La table est partionnée en fragments (\textit{tablets}) qui sont eux-mêmes indexés 
par une racine. 

\vfill

Les fragments pleins subissent un \alert{split}, le nouveau fragment est propagé 
au père.


\end{frame}
 

\begin{frame}
\frametitle{Le routeur}

Le routeur est incorporé au \textit{driver} et maintient une \alert{image} de l'arbre.


\figSlideWithSize{clientImage}{7cm}

Cette image peut être \alert{partielle} et \alert{obsolète}. Elle est \alert{ajustée}
à chaque requête.

\end{frame}
 
\begin{frame}
\frametitle{Exemple d'un ajustement}

Une requête Client transmise sur la base de l'image échoue, car l'arbre a 
beaucoup évolué.

\figSlide{imageadjust}

L'ajustement consiste à remonter/descendre dans l'arbre.

\end{frame}

